\subsection{一致结构的逆像；一致子空间}

设 X 是一个集合，Y 是一致空间，f 是 X 到 Y 的映射。X 上使 f 一致连续的最粗一致结构 U 称为 Y 的一致结构在 f 下的逆像。由第 3 目的命题 4 以及给出交的逆像的公式可知，Y 的一致邻域在 \(g = f \times f\) 下的诸逆像构成 U 的一致邻域基本系。由 U 诱导的拓扑是 Y 的拓扑在 f 下的逆像（第 3 目命题 4 的推论）。

注 若 \(f: \mathbf{X} \to \mathbf{Y}\) 是满射，则 \(\mathbf{Y}\) 的一致邻域是在 \(g\) 之下 \(\mathbf{X}\) 的一致邻域的直接像。

设 \(f\) 是一致空间 \(\mathbf{X}\) 到一致空间 \(\mathbf{X}'\) 的映射；它一致连续的充要条件是：在 \(f\) 之下 \(\mathbf{X}'\) 的一致结构的逆像粗于 \(\mathbf{X}\) 的一致结构。

设 A 是一致空间 X 的子集。X 的一致结构在 A 上诱导的一致结构是后者在典范单射 \(A \rightarrow X\) 下的逆像。由第 3 目命题 4，这等价于下述定义：

定义 3 设 A 是一致空间 X 的子集。A 上以 X 的一致邻域集在 A × A 上的迹为一致邻域集的那个一致结构，称为 X 的一致结构在 A 上诱导的一致结构。

由 A 上诱导的一致结构所诱导的拓扑，与 X 的拓扑在 A 上诱导的拓扑相同；集合 A 连同由 X 的一致结构与拓扑诱导的一致结构与拓扑，称为 X 的一致子空间。

若 A 是一致空间 X 的子集，且 \(f: X \to X'\) 是一致连续映射，则限制 \(f|A\) 是 A 到 \(X'\) 的一致连续映射。若 \(A' \subset X'\) 使得 \(f(X) \subset A'\)，则 X 到一致子空间 \(A'\)（\(X'\) 的）中的、与 f 有同一图像的那个映射仍一致连续（第 3 目命题 4）。

若 \(B \subset A \subset X\)，则 X 的一致子空间 B 与 X 的一致子空间 A 的一致子空间 B 相同（诱导一致结构的可递性；第 3 目命题 5）。

命题 6 设 A 是一致空间 X 的稠密子集。则一致子空间 A 的诸一致邻域在 X × X 中的闭包构成 X 的一致邻域基本系。

A × A 在 X × X 中稠密（第一章 § 4 第 3 目命题 7）。设 V 是 A 的一个开一致邻域；它是 A × A 与 X 的某个开一致邻域 U 的交。我们有 U ⊂ V̄（第一章 § 1 第 6 目命题 5），而这一关系与 V̄ ⊂ U 合在一起，依据 § 1 第 2 目命题 2 的推论 2，便证明了本命题。

